\section{赛前与故障救援：让笔记本变成硬件控制台}
\label{sec:deploy-kvm}

快要上场时，无线连接可能受到现场条件或赛事规定限制。看起来又得抱着便携屏和键盘蹲回地上了？先别急：如果机器人系统与网络服务正常，第~\ref{sec:deploy-wired}~节的\textbf{有线 SSH} 就能解决不少问题。若连网卡都没配好、系统还停在启动界面，或你需要操作 BIOS，则可以准备\textbf{键盘、视频、鼠标（Keyboard, Video, Mouse，KVM）}工具。

这里说的是硬件 KVM：它接收目标电脑的视频输出，并模拟 USB 键鼠，把这些信号送到你的笔记本。它不要求机器人先装好远程桌面服务。和前面的 SSH、RDP 相比，这种独立于目标系统网络服务的访问通常称为\textbf{带外管理（out-of-band management）}。

图~\ref{fig:deploy-kvm-wiring}~先把两种接法放在一起。顺着箭头看：机器人视频送到 KVM，键鼠输入则从 KVM 送回机器人；笔记本接的是 KVM 的控制端，不能省掉机器人侧的线。
\begin{figure}[htbp]
  \centering
  \includegraphics[width=0.95\textwidth]{deployment/kvm_wiring_revised.png}
  \caption{USB KVM 与 IP KVM 的接线关系示意。“笔记本侧一根线”与“机器人侧多条连接”同时成立；供电方式须按具体型号核对。}
  \label{fig:deploy-kvm-wiring}
\end{figure}

\subsubsection{USB KVM：一根线接笔记本，另一侧接机器人}

以 \textbf{Openterface Mini-KVM} 为例，图~\ref{fig:deploy-usb-kvm} 展示设备与配套线缆。笔记本安装厂商的控制端应用，通过一根\emph{支持数据传输的} USB 线接到 KVM 的 Host 端，就可以用自己的屏幕和键鼠操作机器人。机器人端则要把 HDMI 视频输出接到 KVM 的 Input，把机器人 USB 口接到 KVM 的 Target，使其收到模拟键鼠输入。机器人的正常电源仍要保留；这台 Mini-KVM 自身可通过 USB 供电，不代表笔记本也在为机器人供电。

\begin{figure}[htbp]
  \centering
  \includegraphics[width=0.72\textwidth]{deployment/openterface-mini-kvm.jpg}
  \caption{Openterface Mini-KVM 及配套线缆官方照片。橙色长线连接笔记本，短 USB 线与 HDMI 线连接目标电脑；“一根线”描述的是笔记本一侧。}
  \label{fig:deploy-usb-kvm}
\end{figure}

\begin{example}[不依赖网络，操作机载电脑]
先按\href{https://docs.openterface.com/app/kvm/}{官方应用下载页}在笔记本安装对应 Windows、macOS 或 Linux 的 Openterface 控制端。Linux 应选择匹配架构的包，并按其安装说明处理 USB/视频设备权限；机器人端不用安装这款应用。
\begin{enumerate}
  \item 让机器人处于安全测试状态，确认其显示输出接口及所需适配器。在连接处看清 \texttt{Host}、\texttt{Target}、\texttt{Input} 标识，避免只凭 USB-C 外形判断用途。
  \item 用数据 USB 线连接笔记本与 Host；用另一条 USB 线连接机器人与 Target；再把机器人 HDMI 输出接到 KVM 的 HDMI Input。核对两台电脑的供电，不向不明电源接口送电。
  \item 打开笔记本控制端，选择检测到的 KVM 视频设备与控制设备。看到机器人画面后，在应用内点击或输入，打开终端运行 \texttt{hostname}，确认图像和键鼠两条链路都通了。
  \item 在事先约定的维护窗口，可重启\emph{测试机}，观察启动画面并按该主板说明进入 BIOS/UEFI；只查看，不改启动项。由此检查该板在开机阶段能否识别 KVM，不能拿运行中桌面的成功替代这一步。
\end{enumerate}
看得见却点不动时，优先核对 Target USB、数据线与控制权限；完全没有画面时，检查 HDMI 输入、分辨率与输出模式。两者都失败才一起查 Host 端连接和应用设备选择。具体接线对照\href{https://docs.openterface.com/products/minikvm/how-to-connect/}{官方接线指南}。
\end{example}

\begin{note}
多数笔记本的 HDMI 口是\emph{输出}，拿一根 HDMI 线把两台电脑相连，笔记本通常不会变成显示器。普通 HDMI 采集卡可以把画面送入笔记本，但它没有自动提供 USB 键鼠模拟功能；完整控制还需要相应输入链路。USB-C 同样只是接口形状，普通两台电脑之间也不能靠随便一根 USB 线直接互控。
\end{note}

\subsubsection{IP KVM：通过网线进入浏览器控制台}

另一类是 \textbf{IP KVM（KVM over IP）}，例如 \textbf{PiKVM}。图~\ref{fig:deploy-ip-kvm} 中可看到它自己的以太网口：HDMI 和 USB 仍连接机器人，笔记本则通过网络访问\emph{KVM 本身}的网页。机器人自己的网卡就算配置错了，只要 KVM 的网络、供电与视频输入正常，仍有机会从网页控制台进去修复。

\begin{figure}[htbp]
  \centering
  \includegraphics[width=0.78\textwidth]{deployment/pikvm-v4.jpg}
  \caption{PiKVM V4 Mini（左）与 Plus（右）官方照片。网线连接 KVM 的网络端口；其他 RJ45 外形接口可能承担串口或 ATX 功能，必须按标识区分。}
  \label{fig:deploy-ip-kvm}
\end{figure}

\begin{example}[在有线局域网里使用 PiKVM]
以已装好系统的 PiKVM V4 为例，按型号手册接好机器人 HDMI 输出和 USB 键鼠模拟口，并使用规定的独立电源。用网线将 PiKVM 与笔记本接入同一个带 DHCP 的有线调试网络，从 PiKVM 屏幕或路由器租约表读取其地址，例如 \texttt{192.168.50.30}。

在笔记本浏览器打开 \texttt{https://192.168.50.30/}。原厂默认配置的网页账户为 \texttt{admin}、密码为 \texttt{admin}，双因素认证码默认未启用；已配置过的设备应向管理者取得实际凭据，不能靠反复尝试默认密码接管未知设备。自签名证书提示出现时，先通过本地屏幕、租约与接线核对设备身份。

\Needspace{11\baselineskip}
首次登录后先点 Terminal，进入的是 \emph{PiKVM 自己的管理终端}。运行 \texttt{su -}，原厂系统管理员 \texttt{root} 的默认密码为 \texttt{root}，随后逐行运行：
\begin{tbash}
rw
passwd root
kvmd-htpasswd set admin
ro
exit
\end{tbash}
\texttt{rw}/\texttt{ro} 临时切换 PiKVM 的系统分区可写/只读状态；中间两条命令分别交互式设置系统管理员和网页管理员的新密码。退出终端，重新登录网页验证新密码，然后点击 KVM 进入机器人的控制台，打开目标机终端核对 \texttt{hostname}。网页的 KVM 画面是机器人，PiKVM 的 Web Terminal 是 KVM 管理系统，两者不要弄反。接线、出厂凭据与首次配置见\href{https://docs.pikvm.org/v4/}{PiKVM V4 官方快速指南}。

如果想把笔记本用一根网线\emph{直接}接到 PiKVM，而不经过路由器，应先按对应固件说明为它与笔记本配置同网段静态地址；两台设备都在等待 DHCP 时，插线不会凭空出现 DHCP 服务器。先在有 DHCP 的网络完成首次配置，再验证脱网直连，赛前就不用临场猜地址。
\end{example}

IP KVM 的“一根网线”也只描述笔记本侧的访问链路，KVM 到机器人仍需视频、USB 与相应供电。某些型号还有 ATX 开关机控制，但必须另接受支持的主板接口；有 KVM 不等于自动获得机器人总电源控制。板子没有显示输出、供电损坏，或系统卡死后根本不响应 USB，KVM 也不能把它凭空救活。确认具体主板的输出、开机键盘识别与 KVM 固件兼容性，比背产品参数更有用。

\begin{exercise}\label{ex:deploy-kvm}
给“系统正常但无线不能用”“机器人 IP 配错”“停在 BIOS”“目标机完全断电”四种场景分别选择有线 SSH、USB KVM 或 IP KVM，并说明成立条件。再画一张 KVM 接线图，标清笔记本侧和机器人侧各要几条线、谁提供电源；有设备时现场演示“有画面”和“能输入”两项独立验证。
\end{exercise}

\begin{conclusion}
系统和网络服务正常，先用有线 SSH；需要窗口且桌面正常，用 RDP 或 RustDesk；需要查看启动过程、修改网络或绕过失效的桌面服务，考虑硬件 KVM。USB KVM 让笔记本通过 USB 充当控制台，IP KVM 让浏览器通过网络访问控制台。它们都需要目标机具备可用的视频与输入接口；“一根线”不取消机器人侧的接线与供电。
\end{conclusion}

上场前把线缆、转接头、备用地址、登录方式和恢复步骤装进调试包，在本机演练一次，再按当届赛事规则及现场裁判要求确认哪些连接允许使用、何时必须拔除。不能因为改成有线就默认比赛全过程都允许连接；正式启动前还要清点外接线缆，防止运动时拉扯接口。
